\documentclass[border=12pt, varwidth=16cm]{standalone} 
\usepackage[UTF8]{ctex}
\usepackage{amsmath}
\usepackage{amssymb}
\usepackage{color}

\begin{document}

\linespread{1.35}\selectfont 

\noindent
\textbf{残差表示 (Residual Representation)}

\vspace{0.2cm}

\noindent
基于 He et al. (2016) 提出的 \textbf{残差学习假设 (Residual Learning Hypothesis)}：

\vspace{0.2cm}

设 $\mathcal{H}(\mathbf{x})$ 为若干堆叠层 (Stacked Layers) 需要拟合的潜在映射 (Underlying Mapping)，其中 $\mathbf{x}$ 为输入。如果不直接拟合 $\mathcal{H}(\mathbf{x})$，而是让网络去拟合**残差映射 (Residual Mapping)**：
\begin{equation*}
    \mathcal{F}(\mathbf{x}) := \mathcal{H}(\mathbf{x}) - \mathbf{x}
\end{equation*}

\noindent
此时，原始的映射被重构为：
\begin{equation*}
    \mathcal{H}(\mathbf{x}) = \mathcal{F}(\mathbf{x}) + \mathbf{x}
\end{equation*}

\vspace{0.1cm}

\noindent
\textbf{优化的直觉 (Intuition for Optimization):}
\begin{itemize}
    \item 极端的例子：如果最优函数是\textbf{恒等映射 (Identity Mapping)}，即 $\mathcal{H}(\mathbf{x}) = \mathbf{x}$。
    \item 传统方法需要求解器去学习参数以逼近恒等矩阵（这很难）。
    \item 残差方法只需要将多个非线性层的权重推向 \textbf{0}（使得 $\mathcal{F}(\mathbf{x}) \to 0$），这在优化上要容易得多。
\end{itemize}

\end{document}